\section*{September 7, 2026: What is in the treated price distribution's upper tail?}
\textit{Codex draft; follow-up to the price-distribution diagnostics.}

Four school sites account for 82.7\% of the sale value in the treated group's
upper price quarter after closure: Trumbull, Peabody, Humboldt/Duprey, and
Lafayette. These are descriptive concentrations, not estimated closure effects.

Here ``upper quarter'' means sales above the 75th percentile within each
treatment group and year. Pre-closure statistics pool 2008--2012; post-closure
statistics pool 2014--2018. All prices are in 2022
dollars, and each sale has equal weight. Both the percentile thresholds and
the properties above them can change over time.

\begin{center}
{\small\input{../output/upper_tail_table.tex}}
\end{center}

The largest contribution to the widening raw mean gap comes from the
75th--95th percentile band, not the most extreme 5\%. Of the \$58,552 widening
between the pooled pre- and post-closure periods, the 75th--95th band accounts
for \$34,970 (59.7\%), and the top 5\% for \$9,750 (16.7\%). The 50th--75th
band contributes \$12,779 and the bottom half \$1,053. These contributions
sum to the raw mean difference; they do not decompose the controlled regression.

In treated areas, the mean price in the 75th--95th band rises from \$458,987
to \$699,241. Apartment buildings comprise 71.8\% of this band before closure
and 71.7\% afterward, while average building area rises only from 2,816 to
2,887 square feet. The price increase therefore is not accompanied by a broad
shift into apartment buildings or a comparably large increase in building size.
These summaries do not hold every characteristic or location fixed.

The existing property controls include property class (including apartment
buildings), apartment count, building and land area, age, bedrooms, rooms,
full bathrooms, residential type, construction quality, and repair condition.
They reduce the pooled dollar coefficient from \$62,880 to \$56,983, about
9\%. Thus the broad house/apartment distinction is already represented in
the model and does not explain most of its dollar estimate. The controls allow
for observed differences according to the specified formula; they cannot
guarantee constant quality. They also impose common coefficients across years
and sites. Unrecorded renovation, differences within a school area, or changing
market premiums can remain. School-site fixed effects absorb persistent price
levels, not different neighborhood appreciation paths.

Mean--median divergence is useful evidence about the raw distribution,
but a regression of log price estimates average log prices, not the median.
Consequently the raw quantiles alone do not explain the controlled log--dollar
difference. The evidence now locates much of the high-price activity in a
small group of school areas and in a broad upper-price segment. Identifying
the economic cause requires examining their within-area price and property
changes; neighborhood appreciation and unmeasured quality changes remain
hypotheses, not established mechanisms.

{\small Reproduction: Make in the \texttt{home\_price\_distribution} audit
now produces price-band and school-site summaries, each with an automatic
report. The existing source vintage, sample, and regressions are unchanged.
The two new datasets contain 16 band-period rows and 154 site-period rows;
assertions reconstruct group means and upper-quarter contributions. The code is in the current working tree.}
